\documentclass[10pt,a4paper]{amsart}
\usepackage[margin=19mm]{geometry}
\usepackage{amsmath,amssymb,mathtools}
\usepackage[scr=rsfs,cal=euler]{mathalfa}
\usepackage{xfrac}
\usepackage[unicode,hidelinks,bookmarks=false]{hyperref}
\newcommand{\ie}{i.e.\ }
\newcommand{\cf}{cf.\ }
\newcommand{\resp}{respectively\ }
\newcommand{\Subsection}{\subsection}
\providecommand{\nobreakdash}{\nobreak\hbox{-}}
\setlength{\emergencystretch}{2em}
\multlinegap0pt
\usepackage{amsthm}
\usepackage{mathtools}
\usepackage{amsmath,amssymb,amsfonts}
\usepackage{graphicx}
\usepackage{bm}
\usepackage{enumitem}
\usepackage{mathrsfs}
\usepackage{xcolor}
\theoremstyle{plain}
\newtheorem{theorem}{Theorem}[section]
\newtheorem{lemma}[theorem]{Lemma}
\newtheorem{proposition}[theorem]{Proposition}
\newtheorem{corollary}[theorem]{Corollary}
\newtheorem{definition}[theorem]{Definition}
\newtheorem{remark}[theorem]{Remark}
\newtheorem*{acknowledgment}{Acknowledgments}
\newtheorem{Assumption}{H\!\!} % Custom label for Assumption
\numberwithin{equation}{section}
\title[Smoluchowski-Kramers approximation with Levy noise]{Smoluchowski-Kramers Approximation with Levy Noise in the Meyer-Zheng Topology: tightness of the slow and transformed families in the Meyer-Zheng topology (Lemma 3.5)}
\author[X. Liu, Q. Jiang, W. Wang]{Xueru Liu, Qianqian Jiang, Wei Wang}
\date{Source: arXiv:2609.29205v1 (CC BY 4.0)}
\begin{document}
\maketitle
\vspace{1em}
\noindent\textit{Source and licence.} Smoluchowski-Kramers approximation with Levy noise in the Meyer-Zheng topology. Version arXiv:2609.29205v1 (CC BY 4.0), distributed under the Creative Commons Attribution 4.0 International licence, http://creativecommons.org/licenses/by/4.0/, as recorded in the arXiv OAI licence metadata for this exact version and shown on the abstract page https://arxiv.org/abs/2609.29205v1. This item is an editorial re-presentation of the paper's tightness lemma for the slow variable and the transformed variable in the Meyer-Zheng topology, printed as Lemma 3.5. A full rights notice is delivered at licenses/arxiv-2609.29205-CC-BY-4.0.txt.
\noindent\textit{Editorial scope.} Counted unit: the paper's tightness lemma for the slow variable and the transformed variable in the Meyer-Zheng topology, printed as Lemma 3.5 (source0.tex lines the retained context spans are the paper's dataset and assumption block, its tightness and projection lemmas, and the conditional-variation lemma with its proof; the counted statement is at lines 1201--1206 and the counted proof at lines 1208--1337), delivered together with the paper's complete original proof of it as the single primary proof body. No mathematics is written, repaired or re-derived: statement and proof are the paper's own text line for line, in the paper's own notation. Required context retained uncounted, because the counted proof is the paper's own assembly over it: the definitions and standing assumptions H1 to H3, the tightness and projection lemmas of the paper's Section 2, the energy estimates, and the conditional-variation lemma with its complete 130-line proof on which the counted argument depends. Every cross-reference inside the retained spans resolves inside the delivered document. Printed numbers are the paper's own: the wrapper restores the paper's counters before the retained material, so the counted statement prints as Lemma 3.5, the paper's own number for it, and the numbered displays and theorem-like items of the retained context carry the wrapper's own sequential numbering because the paper's counters are restored immediately before the counted statement. The retained context is delivered so that the counted proof can be checked line by line; its own printed numbers are the wrapper's sequential numbers rather than the paper's, which is disclosed here. Omitted from this package and not redistributed: the paper's preamble and front matter as such (of which only the title and the author names are reproduced in the wrapper); the sections, lemmas, theorems and examples that lie outside the retained spans; and the paper's bibliography with its bib data file. No citation command occurs in any retained span, so the delivered document carries no bibliography and no reference or citation command of the delivered text was rewritten or adapted. Printed slips kept literal, among them the paper's own wording and its own choice of symbols. Layout and class: the paper's own class is the standard LaTeX article class at 11pt with section-relative equation numbering, which is not a publisher class; the delivered wrapper is the lane's pinned amsart editorial wrapper at 10pt on A4, to which this package adds the paper's own theorem-like declarations with their shared section-relative counter and the shorthand macros its retained text uses, all copied verbatim from the paper's source. No publisher class file, no publisher style file and no upstream script is redistributed. Assets: no retained span of this package contains a figure environment or a graphics-inclusion command, no image or artwork file exists in or is redistributed with this package, and the manifest and the delivery evidence both carry an empty asset-dependency list. A full rights notice is delivered at licenses/arxiv-2609.29205-CC-BY-4.0.txt.
\setcounter{section}{1}
\begin{equation}\label{e:SLE-11}
\left\{
\begin{aligned}
du^\epsilon(t) &=v^\epsilon(t)\,dt,\\
\epsilon\,dv^\epsilon(t) &=\Delta u^\epsilon(t)\,dt-\gamma(u^\epsilon(t))v^\epsilon(t)\,dt+F(u^\epsilon(t))\,dt+\sigma(u^\epsilon(t))\,dW^Q(t)\\
&\quad +\int_{Z_1}G(u^\epsilon(t-),z)\,\tilde{\mathcal N}(dt,dz) +\int_{\mathbb R^m\setminus Z_1} G(u^\epsilon(t-),z)\,\mathcal N(dt,dz),\\
u^\epsilon(t)|_{\partial\mathcal O}&=0,\qquad u^\epsilon(0)=u_0,\qquad v^\epsilon(0)=v_0 ,
\end{aligned}
\right.
\end{equation}

\begin{Assumption}\label{H1}
The measure \(\nu\) is a L\'evy measure on \(\mathbb R^m\), 
\[
\nu(\{0\})=0, \qquad \int_{\mathbb R^m}(1\wedge |z|^2)\,\nu(dz)<\infty .
\]
Moreover, \(\nu\) has finite second moment, that is \(\int_{\mathbb R^m}|z|^2\,\nu(dz)<\infty .\)
\end{Assumption}

\begin{Assumption}\label{H2}
The mappings \(\sigma:H\to L_2(H^Q,H),~F:H\to H,~G:H\times\mathbb R^m\to H\)
are measurable. Moreover, there exist constants \(L_\sigma,L_F,L_G>0\) such that, for all
\(u,u_1,u_2\in H\) and \(z\in\mathbb R^m\),
\begin{align*}
\|\sigma(u_1)-\sigma(u_2)\|_{L_2(H^Q,H)}&\le L_\sigma\|u_1-u_2\|_H,&\|\sigma(u)\|_{L_2(H^Q,H)}&\le L_\sigma(1+\|u\|_H),\\
\|F(u_1)-F(u_2)\|_H&\le L_F\|u_1-u_2\|_H,&\|F(u)\|_H&\le L_F(1+\|u\|_H),\\
\|G(u_1,z)-G(u_2,z)\|_H^2&\le L_G\|u_1-u_2\|_H^2|z|^2,&\|G(u,z)\|_H^2&\le L_G(1+\|u\|_H)^2|z|^2 .
\end{align*}
\end{Assumption}

\begin{Assumption}\label{H3}
The damping coefficient \(\gamma:\mathbb R\to\mathbb R\) belongs to
\(C_b^1(\mathbb R)\) and satisfies
\[
0<\gamma_0\le \gamma(r)\le \gamma_1,\qquad|\gamma'(r)|\le \gamma_1,\qquad r\in\mathbb R,
\]
for some constants \(\gamma_0,\gamma_1>0\). For \(u\in H\),   \(\gamma(u)(x):=\gamma(u(x)).\)
\end{Assumption}

\begin{theorem}[Meyer--Zheng tightness criterion]\label{thm:MZ-tightness}
Let \(\{X^n\}_{n\ge1}\) be a sequence of real-valued adapted c\`adl\`ag
processes on \([0,T]\). Suppose that each \(X^n\) is a quasimartingale and
\[
\sup_{n\ge1}{\rm CV}_T(X^n)<\infty.
\]
Then the laws of \(\{X^n\}_{n\ge1}\) are tight in
\(D([0,T];\mathbb R)\) endowed with the Meyer--Zheng topology.
\end{theorem}

\begin{lemma}[Projection criterion for Meyer--Zheng tightness]
\label{lem:MZ-projection-tightness}
Let \(\{Y^n\}_{n\ge1}\) be a sequence of \(H^{-1}\)-valued adapted c\`adl\`ag processes. Assume that, for every \(k\ge1\),
\[
\sup_{n\ge1}{\rm CV}_T\bigl(\langle Y^n,e_k\rangle_{H^{-1},H^1}\bigr)<\infty,
\]
and that, for every \(\eta>0\),
\[
\lim_{N\to\infty}\sup_{n\ge1}\mathbb P\left(\int_0^T\|(I-P_N)Y^n(t)\|_{H^{-1}}\,dt>\eta\right)=0.
\]
Then the laws of \(\{Y^n\}_{n\ge1}\) are tight in \(D([0,T];H^{-1})\) endowed with the Meyer--Zheng topology.
\end{lemma}

\begin{proof}
For each fixed \(N\in\mathbb N\), the process \(P_NY^n\) is
finite-dimensional. By Theorem~\ref{thm:MZ-tightness}, applied to each
coordinate
\[
\langle Y^n,e_k\rangle_{H^{-1},H^1},\qquad 1\le k\le N,
\]
the laws of \(\{P_NY^n\}_{n\ge1}\) are tight in \(D([0,T];P_NH^{-1})\) endowed with the Meyer--Zheng topology.
Moreover,
\[
d_{\rm MZ}(Y^n,P_NY^n)\le\int_0^T\|(I-P_N)Y^n(t)\|_{H^{-1}}\,dt .
\]
The assumed tail estimate therefore implies that \(Y^n-P_NY^n\) is negligible in the Meyer--Zheng topology, uniformly in \(n\), as
\(N\to\infty\). Combining the finite-dimensional tightness with this tail estimate yields the desired tightness of \(\{Y^n\}_{n\ge1}\) in
\(D([0,T];H^{-1})\).
\end{proof}

\begin{equation}\label{a-priori-v}
\sup_{t\in[0,T]}
\Big(
\epsilon\,\mathbb E\|v^{\epsilon}(t)\|_H^2
+ \mathbb E\|\nabla u^{\epsilon}(t)\|_H^2
\Big)
\le C_T .
\end{equation}

\begin{equation}\label{lemma-g-2}
\mathbb{E}\sup_{t\in [0,T]}\|\eta^\epsilon(t)\|_H^2  +\mathbb{E}\int_0^T \|\nabla \eta^{\epsilon}(s)\|_{H}^2\,ds
\leq \gamma_1^2 C_T.
\end{equation}

\begin{align}\label{lemma-g-3}
\eta^\epsilon(t) + \epsilon   v^\epsilon(t) 
&= g(u_0) + \epsilon v_0 
+ \int_0^t \nabla \cdot \left( B(\eta^\epsilon(s)) \nabla \eta^\epsilon(s) \right)\,ds 
+ \int_0^t \bar{F}(\eta^\epsilon(s))\,ds \nonumber \\
&\quad + \int_0^t \sigma_g(\eta^\epsilon(s))\,dW^Q(s)
+ \int_0^t \int_{Z_1} G_g(\eta^\epsilon(s),z)\,\tilde{{\mathcal N}}(ds,dz) \nonumber \\
&\quad + \int_0^t \int_{\mathbb{R}^m \setminus Z_1} G_g(\eta^\epsilon(s),z)\,{\mathcal N}(ds,dz).
\end{align}

\begin{equation}\label{lemma-Br}
\frac{1}{\gamma_1} \leq B(r) \leq \frac{1}{\gamma_0}, \qquad \forall r \in \mathbb{R},
\end{equation}

\begin{lemma}\label{lem:conditional-variation-Yk}
Assume that Assumptions~\ref{H1}--\ref{H3} hold. Then, for every fixed \(k\ge1\),
\[
\sup_{0<\epsilon\le1}{\rm CV}_{T}(Y_k^\epsilon)<\infty .
\]
Consequently, \(\{Y_k^\epsilon\}_{0<\epsilon\le1}\) is tight in \(D([0,T];\mathbb R)\) endowed with the Meyer--Zheng topology.
\end{lemma}

\begin{proof}
By \eqref{lemma-g-3}, for every fixed \(k\ge1\),
\[
\begin{aligned}
Y_k^\epsilon(t)={}&\langle g(u_0)+\epsilon v_0,e_k\rangle_H +\int_0^t\left\langle\nabla\cdot\left(B(\eta^\epsilon(s))\nabla\eta^\epsilon(s)\right),e_k
\right\rangle_{H^{-1},H^1} ds  \\
&+\int_0^t\langle \bar F(\eta^\epsilon(s)),e_k\rangle_H\,ds   +\int_0^t\langle \sigma_g(\eta^\epsilon(s))\,dW^Q(s),e_k\rangle_H  \\
&+\int_0^t\int_{Z_1}\langle G_g(\eta^\epsilon(s-),z),e_k\rangle_H\,\tilde{\mathcal N}(ds,dz)    \\
&+\int_0^t\int_{\mathbb R^m\setminus Z_1}\langle G_g(\eta^\epsilon(s-),z),e_k\rangle_H\,\mathcal N(ds,dz).
\end{aligned}
\]
Writing the large-jump integral as the sum of its compensated part and its
compensator, we obtain
\[
Y_k^\epsilon(t)=Y_k^\epsilon(0)+\int_0^t a_k^\epsilon(s)\,ds+M_k^\epsilon(t),
\]
where
\[
\begin{aligned}
a_k^\epsilon(s):={}&\left\langle\nabla\cdot\left(B(\eta^\epsilon(s))\nabla\eta^\epsilon(s)\right),e_k\right\rangle_{H^{-1},H^1} +\langle \bar F(\eta^\epsilon(s)),e_k\rangle_H   \\
&+\int_{\mathbb R^m\setminus Z_1}\langle G_g(\eta^\epsilon(s),z),e_k\rangle_H\,\nu(dz),
\end{aligned}
\]
and
\[
\begin{aligned}
M_k^\epsilon(t):={}&\int_0^t\langle \sigma_g(\eta^\epsilon(s))\,dW^Q(s),e_k\rangle_H    +\int_0^t\int_{Z_1}
\langle G_g(\eta^\epsilon(s-),z),e_k\rangle_H\, \tilde{\mathcal N}(ds,dz)                                            \\
&+\int_0^t\int_{\mathbb R^m\setminus Z_1}\langle G_g(\eta^\epsilon(s-),z),e_k\rangle_H\,\tilde{\mathcal N}(ds,dz).
\end{aligned}
\]
By Assumptions~\ref{H1}--\ref{H3} and the uniform estimate \eqref{lemma-g-2}, \(M_k^\epsilon\) is a square-integrable real-valued
martingale.
Let \(\pi: 0=t_0<t_1<\cdots<t_n=T \) be an arbitrary partition of \([0,T]\). Since \(M_k^\epsilon\) is a martingale,
\(\mathbb E\left[ M_k^\epsilon(t_{i+1})-M_k^\epsilon(t_i)\mid \mathcal F_{t_i}\right]=0 .\)
Hence
\[
\begin{aligned}
\mathbb E\sum_{i=0}^{n-1}\left|\mathbb E\left[Y_k^\epsilon(t_{i+1})-Y_k^\epsilon(t_i)\mid \mathcal F_{t_i}\right]\right|  
 &= \mathbb E\sum_{i=0}^{n-1}\left|\mathbb E\left[\int_{t_i}^{t_{i+1}} a_k^\epsilon(s)\,ds\mid \mathcal F_{t_i}\right]\right|          \\
&\quad \le \mathbb E\sum_{i=0}^{n-1}\mathbb E\left[\int_{t_i}^{t_{i+1}} |a_k^\epsilon(s)|\,ds\mid \mathcal F_{t_i}\right]   \\
&\quad = \mathbb E\int_0^T |a_k^\epsilon(s)|\,ds .
\end{aligned}
\]
It remains to estimate the three terms in \(a_k^\epsilon\). First, using the definition of the distributional divergence and
\eqref{lemma-Br}, we have
\[
\begin{aligned}
\left|\left\langle\nabla\cdot\left(B(\eta^\epsilon)\nabla\eta^\epsilon\right),e_k\right\rangle_{H^{-1},H^1}\right|
&=\left|-\int_{\mathcal O}B(\eta^\epsilon)\nabla\eta^\epsilon\cdot\nabla e_k\,dx\right|                                                
\le\frac1{\gamma_0}\|\nabla\eta^\epsilon\|_H\|\nabla e_k\|_H .\end{aligned}
\]
Therefore, by \eqref{lemma-g-2},
\[
\begin{aligned}
 \mathbb E\int_0^T\left|\left\langle\nabla\cdot\left(B(\eta^\epsilon(s))\nabla\eta^\epsilon(s)\right),e_k\right\rangle_{H^{-1},H^1}\right|ds                                          
 &  \le\frac1{\gamma_0}\|\nabla e_k\|_H\int_0^T\left(\mathbb E\|\nabla\eta^\epsilon(s)\|_H^2\right)^{1/2}ds\\
&\leq\frac{\sqrt{T}}{\gamma_0}|\nabla e_k|_H\left(\mathbb E\int_0^T|\nabla\eta^\epsilon(s)|_H^2,ds\right)^{1/2} \le C_{T,k}.
\end{aligned}
\]
Second, by Assumptions~\ref{H2}--\ref{H3} and the Lipschitz continuity of \(g^{-1}\), \(\|\bar F(\eta)\|_H=\|F(g^{-1}(\eta))\|_H\le C(1+\|\eta\|_H).\)
Thus, using again \eqref{lemma-g-2},
\[
\begin{aligned}
\mathbb E\int_0^T |\langle \bar F(\eta^\epsilon(s)),e_k\rangle_H|\,ds\le\int_0^T\mathbb E\|\bar F(\eta^\epsilon(s))\|_H\,ds        
\le C\int_0^T\left(1+\mathbb E\|\eta^\epsilon(s)\|_H\right)ds \le C_T .
\end{aligned}
\]
Third, by Assumptions~\ref{H1}--\ref{H3},  the Lipschitz continuity of \(g^{-1}\) and  \eqref{lemma-g-2},
%\[
%\begin{aligned}
%\int_{\mathbb R^m\setminus Z_1}
%\|G_g(\eta,z)\|_H\,\nu(dz)
%&=
%\int_{\mathbb R^m\setminus Z_1}
%\|G(g^{-1}(\eta),z)\|_H\,\nu(dz)                       \\
%&\le
%C(1+\|g^{-1}(\eta)\|_H)
%\int_{\mathbb R^m\setminus Z_1}|z|\,\nu(dz)             \\
%&\le
%C(1+\|\eta\|_H).
%\end{aligned}
%\]
%Indeed, by Assumption~\ref{H1},
%\[
%\int_{\mathbb R^m\setminus Z_1}|z|\,\nu(dz)
%\le
%\left(
%\int_{\mathbb R^m\setminus Z_1}|z|^2\,\nu(dz)
%\right)^{1/2}
%\left(
%\nu(\mathbb R^m\setminus Z_1)
%\right)^{1/2}
%<\infty .
%\]
%Hence, by \eqref{lemma-g-2},
\[
\begin{aligned}
\mathbb E\int_0^T \int_{\mathbb R^m\setminus Z_1}|\langle G_g(\eta^\epsilon(s),z),e_k\rangle_H|\,\nu(dz)\,ds                                          
&\le\mathbb E\int_0^T\int_{\mathbb R^m\setminus Z_1}\|G_g(\eta^\epsilon(s),z)\|_H\,\nu(dz)\,ds              \\
&\le C\int_0^T\left(1+\mathbb E\|\eta^\epsilon(s)\|_H\right)ds \le C_T .
\end{aligned}
\]
Combining the preceding estimates gives
\[
\sup_{\pi}\mathbb E\sum_{i=0}^{n-1}\left|\mathbb E\left[ Y_k^\epsilon(t_{i+1})-Y_k^\epsilon(t_i)\mid \mathcal F_{t_i}\right]\right|\le C_{T,k},
\]
where \(C_{T,k}\) is independent of \(0<\epsilon\le1\).

Since \(Y_k^\epsilon(T)=\langle \eta^\epsilon(T),e_k\rangle_H+\epsilon\langle v^\epsilon(T),e_k\rangle_H ,\)
and \(\|e_k\|_H=1\),  by \eqref{lemma-g-2} and \eqref{a-priori-v}, we obtain
\[
\begin{aligned}
\mathbb E|Y_k^\epsilon(T)| \le \mathbb E\|\eta^\epsilon(T)\|_H+ \epsilon\,\mathbb E\|v^\epsilon(T)\|_H            
&\le C_T+\epsilon\left(\mathbb E\|v^\epsilon(T)\|_H^2\right)^{1/2}               \\
&\le C_T+C_T\sqrt{\epsilon}\le C_T , \qquad 0<\epsilon\le1 .
\end{aligned}
\]
Therefore,
\[
\sup_{0<\epsilon\le1}{\rm CV}_{T}(Y_k^\epsilon) \le C_{T,k}<\infty .
\]
By Theorem~\ref{thm:MZ-tightness}, the family \(\{Y_k^\epsilon\}_{0<\epsilon\le1}\) is tight in  \(D([0,T];\mathbb R)\) endowed with the Meyer--Zheng topology.
The proof is complete.
\end{proof}

\begin{lemma}\label{thm:main3}
Assume that Assumptions~\ref{H1}--\ref{H3} hold. Let \((u^\epsilon,v^\epsilon)\) be the solution to~\eqref{e:SLE-11} with
initial condition \((u_0,v_0)\in H^1\times H .\) Then, for every fixed \(T>0\), the families of laws
\(\left\{\mathcal L(Y^\epsilon)\right\}_{0<\epsilon\le1} ~\text{and}~ \left\{\mathcal L(\eta^\epsilon)\right\}_{0<\epsilon\le1}\)
are tight in \(D([0,T];H^{-1})\) endowed with the Meyer--Zheng topology.
\end{lemma}

\begin{proof}
We first prove the tightness of \(\{Y^\epsilon\}_{0<\epsilon\le1}\) in \(D([0,T];H^{-1})\) endowed with the Meyer--Zheng topology.
For each fixed \(k\ge1\), by Lemma~\ref{lem:conditional-variation-Yk}, \(\sup_{0<\epsilon\le1}{\rm CV}_{T}(Y_k^\epsilon)<\infty .\)
Hence the coordinate condition in Lemma~\ref{lem:MZ-projection-tightness} is satisfied

Since \(Y^\epsilon=\eta^\epsilon+\epsilon v^\epsilon,\) we have \((I-P_N)Y^\epsilon=(I-P_N)\eta^\epsilon+(I-P_N)(\epsilon v^\epsilon).\)
Hence,
\[
\|(I-P_N)Y^\epsilon(t)\|_{H^{-1}}^2 \le 2\|(I-P_N)\eta^\epsilon(t)\|_{H^{-1}}^2 + 2\|(I-P_N)(\epsilon v^\epsilon(t))\|_{H^{-1}}^2 .
\]
Firstly,  by \(\eta^\epsilon(t)=\sum_{k=1}^\infty \eta_k^\epsilon(t)e_k,\)
then
\(
\|(I-P_N)\eta^\epsilon(t)\|_{H^{-1}}^2 = \sum_{k=N+1}^\infty \alpha_k^{-1}|\eta_k^\epsilon(t)|^2 .
\)
For \(k\ge N+1\), since \(\alpha_k\ge \alpha_{N+1}\), \(\alpha_k^{-1} \le \alpha_{N+1}^{-2}\alpha_k . \) Hence
\[
\begin{aligned}
\|(I-P_N)\eta^\epsilon(t)\|_{H^{-1}}^2 \le\alpha_{N+1}^{-2}\sum_{k=N+1}^\infty\alpha_k|\eta_k^\epsilon(t)|^2 
 \le \alpha_{N+1}^{-2}\|\eta^\epsilon(t)\|_{H^1}^2 .
\end{aligned}
\]
Consequently, by \eqref{lemma-g-2},
\[
\begin{aligned}
\mathbb E\int_0^T \|(I-P_N)\eta^\epsilon(t)\|_{H^{-1}}^2\,dt
 \le \alpha_{N+1}^{-2}\mathbb E\int_0^T\|\eta^\epsilon(t)\|_{H^1}^2\,dt  \le C_T\alpha_{N+1}^{-2}.
\end{aligned}
\]
Next, by \(v^\epsilon(t)=\sum_{k=1}^\infty v_k^\epsilon(t)e_k,\) then,
\(
\|(I-P_N)(\epsilon v^\epsilon(t))\|_{H^{-1}}^2=\sum_{k=N+1}^\infty\alpha_k^{-1}\epsilon^2|v_k^\epsilon(t)|^2 .
\)
Since \(\alpha_k^{-1}\le \alpha_{N+1}^{-1}\) for \(k\ge N+1\), we obtain
\[
\begin{aligned}
\|(I-P_N)(\epsilon v^\epsilon(t))\|_{H^{-1}}^2 \le\alpha_{N+1}^{-1}\epsilon^2\sum_{k=N+1}^\infty |v_k^\epsilon(t)|^2    
 \le \alpha_{N+1}^{-1}\epsilon^2\|v^\epsilon(t)\|_H^2 .
\end{aligned}
\]
Using \eqref{a-priori-v}, we have
\[
\begin{aligned}
\mathbb E\int_0^T
\|(I-P_N)(\epsilon v^\epsilon(t))\|_{H^{-1}}^2\,dt
&\le
\alpha_{N+1}^{-1}
\epsilon^2
\int_0^T
\mathbb E\|v^\epsilon(t)\|_H^2\,dt        \\
&\le
C_T\alpha_{N+1}^{-1}\epsilon    \le
C_T\alpha_{N+1}^{-1},
\qquad 0<\epsilon\le1 .
\end{aligned}
\]
Combining the preceding estimates yields
\[
\mathbb E\int_0^T
\|(I-P_N)Y^\epsilon(t)\|_{H^{-1}}^2\,dt
\le
C_T\alpha_{N+1}^{-2}
+
C_T\alpha_{N+1}^{-1}.
\]
Since \(\alpha_{N+1}\to\infty\), it follows that
\[
\lim_{N\to\infty}
\sup_{0<\epsilon\le1}
\mathbb E\int_0^T
\|(I-P_N)Y^\epsilon(t)\|_{H^{-1}}^2\,dt
=0.
\]
In particular, for every \(\delta>0\),
\[
\begin{aligned}
&\sup_{0<\epsilon\le1}
\mathbb P
\left(
\int_0^T
\|(I-P_N)Y^\epsilon(t)\|_{H^{-1}}\,dt>\delta
\right)                                                     \\
&\quad\le
\frac{T^{1/2}}{\delta}
\left(
\sup_{0<\epsilon\le1}
\mathbb E\int_0^T
\|(I-P_N)Y^\epsilon(t)\|_{H^{-1}}^2\,dt
\right)^{1/2}
\longrightarrow 0,~~ as ~N\to\infty. 
\end{aligned}
\]
Thus the  tail condition in
Lemma~\ref{lem:MZ-projection-tightness} holds. Therefore, \(\{Y^\epsilon\}_{0<\epsilon\le1}\)
is tight in \(D([0,T];H^{-1})\) endowed with the Meyer--Zheng topology.

It remains to prove the tightness of
\(\eta^\epsilon\). By definition, \(Y^\epsilon(t)-\eta^\epsilon(t)=\epsilon v^\epsilon(t).\)
Using the continuous embedding \(H\subset H^{-1}\) and \eqref{a-priori-v}, we obtain
\[
\begin{aligned}
\mathbb E\int_0^T\|Y^\epsilon(t)-\eta^\epsilon(t)\|_{H^{-1}}^2\,dt
=\mathbb E\int_0^T\|\epsilon v^\epsilon(t)\|_{H^{-1}}^2\,dt                  
\le C\epsilon^2\int_0^T\mathbb E\|v^\epsilon(t)\|_H^2\,dt                
\le C_T\epsilon .
\end{aligned}
\]
Hence \(Y^\epsilon-\eta^\epsilon \longrightarrow 0 ~ \text{in }L^2(\Omega\times[0,T];H^{-1}).\)
Moreover,
\[
d_{\rm MZ}(Y^\epsilon,\eta^\epsilon) \le \int_0^T \|Y^\epsilon(t)-\eta^\epsilon(t)\|_{H^{-1}}\,dt .
\]
Therefore,
\[
d_{\rm MZ}(Y^\epsilon,\eta^\epsilon) \longrightarrow 0 \qquad \text{in probability}.
\]
Since \(\{Y^\epsilon\}_{0<\epsilon\le1}\) is tight in \(D([0,T];H^{-1})\) endowed with the Meyer--Zheng topology, for every
sequence \(\epsilon_n\to0\), there exists a subsequence, still denoted by \(\epsilon_n\), such that \(Y^{\epsilon_n}\Longrightarrow Y\)
in \(D([0,T];H^{-1})\) endowed with the Meyer--Zheng topology. Since
\[
d_{\rm MZ}(Y^{\epsilon_n},\eta^{\epsilon_n})\longrightarrow 0\qquad \text{in probability},
\]
it follows from Slutsky's theorem that \(\eta^{\epsilon_n}\Longrightarrow Y\) in \(D([0,T];H^{-1})\) endowed with the Meyer--Zheng topology. 
Hence every sequence \(\{\eta^{\epsilon_n}\}_{n\ge1}\) admits a weakly convergent subsequence. Therefore, \(\{\eta^\epsilon\}_{0<\epsilon\le1}\) 
is tight in \(D([0,T];H^{-1})\) endowed with the Meyer--Zheng topology.
 


 
\end{proof}

\end{document}
